% Cover Letter — Quantbot | Quantitative Researcher Internship - 2027 | Hong Kong
\documentclass[a4paper,11pt]{article}
\usepackage[empty]{fullpage}
\usepackage{geometry}
\usepackage[hidelinks]{hyperref}
\usepackage{parskip}
\usepackage[T1]{fontenc}

\geometry{left=2.5cm, top=2cm, right=2.5cm, bottom=2cm}
\pagestyle{empty}
\setlength{\parindent}{0pt}
\setlength{\parskip}{10pt}

\begin{document}

\begin{flushright}
\textbf{Ojas Srivastava} \\
+91-7424978046 \\
\href{mailto:srivastavaojas454@gmail.com}{srivastavaojas454@gmail.com} \\
\href{https://github.com/Ojas-Srivastava05}{GitHub} \textbar{}
\href{https://linkedin.com/in/ojas-srivastava05}{LinkedIn} \\
Surat, India \textbar{} Available for 10-week internship in Hong Kong (Jun--Aug 2027) \\
\today
\end{flushright}

\vspace{4mm}
\textbf{Hiring Team} \\
Quantbot Technologies \\
Quantitative Researcher Internship --- 2027 \\
Hong Kong

\vspace{2mm}
\textbf{Re: Application for Quantitative Researcher Internship - 2027 (Hong Kong)}

Dear Hiring Team,

I am applying for the Quantitative Researcher Internship (2027) in Hong Kong. I am a penultimate-year B.Tech Artificial Intelligence student at SVNIT Surat (CGPA 9.20/10, graduating May 2028), with coursework in probability and statistics, linear algebra, and machine learning, and strong programming skills in Python and C++. I am drawn to Quantbot's systematic, research-first approach---large-scale data, rigorous testing, and modern statistical/ML methods---and I want first-hand experience designing, validating, and refining quantitative ideas with a dedicated mentor over the 10-week program.

My strongest research-style project is LogiFlow (Google Solution Challenge 2026 Global Top 106). I treated delay prediction as an open-ended research problem: formulated the hypothesis, engineered features from timetable, route, and history on 15,650 train-day records, and built Gradient Boosting models with cross-validation (MAE 22.7 minutes; 81\% of predictions within 30 minutes). I iterated on experimental design, built a 100/100 automated validation suite, and developed ranking/explainability tooling---practice moving from question to model to measured outcome, then presenting results clearly.

I am comfortable analyzing large structured datasets in Python, writing clean research code, and communicating findings. Competitive programming (LeetCode Knight, peak 2048; Codeforces Specialist, 1419) keeps my problem-solving sharp when the first approach fails. Separately, I have followed Indian equities and mutual funds for over three years and write small Python tools for portfolio logging and screening---finance knowledge is developing, not assumed. I would be excited to contribute to alpha-oriented research workflows at Quantbot, learn how high-performance computing accelerates experiments, and grow toward a full-time research path if the internship is successful.

Thank you for your consideration. I would welcome the opportunity to discuss how I can contribute.

Sincerely, \\
\vspace{8mm}
\textbf{Ojas Srivastava}

\end{document}
